\chapter{Voids and tessellations}\label{ch:voids}
\chapterstatus{written}{2026-11-26}

The Delaunay tessellation of the particles gives density and velocity fields without a grid (DTFE) and the Voronoi densities of
watershed void finders. \fml{} has its own MPI-parallel Delaunay triangulation and periodic Voronoi tessellation in 2D and 3D, with
exact predicates and results that do not depend on the number of tasks, which can follow moving points (the building block of a
future moving-mesh hydro solver); the older tessellation uses CGAL. ZOBOV, the ORIGAMI halos and the DTFE run on either, the
native one by default; CGAL is optional. \fml{} has four void finders sharing one
catalogue format: spherical, voxel watershed (Revolver-like) and ZOBOV, with optional zone merging.

\begin{modulebox}{Triangulation, VoidFinders}
\nsp{FML::DELAUNAY} (header-only, no dependencies):
\file{Triangulation/Predicates.h}, \file{Delaunay.h} (\cls{Delaunay}), \file{PeriodicVoronoi.h} (\cls{PeriodicVoronoi}).
\nsp{FML::TRIANGULATION}: \file{Triangulation/VoronoiTessellation.h} (\cls{PeriodicVoronoiTessellation}, native),
\file{Triangulation/Watershed.h} (\code{WatershedOnGraph}), \file{Triangulation/PeriodicDelaunay.h} (\cls{MPIPeriodicDelaunay},
CGAL), \file{Triangulation/DTFE.h} (\cls{DTFE}, either), \file{Triangulation/WatershedBinning.h}. \file{VoidFinders/Void.h}
(\cls{Void}), \file{SphericalVoidFinder.h}, \file{VoxelVoidFinder.h} (also \code{VoxelWatershed}), \file{ZobovVoidFinder.h},
\file{ZoneMerging.h}.\par
Documentation: \file{Triangulation/DELAUNAY.md}, \file{Triangulation/README.md}, \file{VoidFinders/README.md}. Tests: \file{Triangulation/tests}
(\file{test_delaunay.cpp}, \file{test_delaunay_mpi.cpp}, \file{compare_cgal.cpp}, \file{compare_library.cpp}),
\file{tests/test_watershed_native.cpp}, \file{tests/test_origami_halos.cpp}, \file{tests/test_dtfe_native.cpp}, \file{VoidFinders/tests} (\file{test_spherical.cpp}),
\file{tests/test_delaunay_domain.cpp}, \file{tests/test_dtfe_pencil.cpp}, \file{tests/test_voids_pencil.cpp}. CGAL \cite{CGAL}
(\code{USE_CGAL}) is needed only for \cls{MPIPeriodicDelaunay} and the options that use it.
\end{modulebox}

\section{FML's own Delaunay triangulation}\label{sec:voids-native}

\paragraph*{Predicates.} Everything rests on two exact predicates, \code{orient} (the sign of
$\det[\bm p_1-\bm p_0,\dots,\bm p_D-\bm p_0]$) and \code{insphere} (+1 if $\bm q$ is inside the circumsphere of a positively
oriented simplex), with CGAL's conventions. The determinant is evaluated in floating point with Shewchuk's a priori error bound
\cite{Shewchuk1997} (times a safety factor 4); when the sign is uncertain it is recomputed exactly with floating-point expansions
of the input coordinates, with a fixed-size stack version when all coordinate differences are exact doubles (lattices, quantized
coordinates). The result is exact for any double input barring overflow.

\paragraph*{Triangulation.} \cls{Delaunay}$\langle D\rangle$ ($D=2,3$) uses incremental insertion \cite{Bowyer1981,Watson1981}, placing the points in Peano--Hilbert order into a super
simplex (Bowyer--Watson): a visibility walk from the last new simplex locates the point, a breadth-first search finds the cavity of
simplices whose circumspheres contain it, and the cavity is replaced by the simplices joining its boundary facets to the point.
Cospherical points (lattices) are resolved by the symbolic perturbation of Devillers \& Teillaud \cite{DevillersTeillaud2011},
as in CGAL: the result is the Delaunay triangulation of the points perturbed in the order of a per-vertex key, a valid and unique
Delaunay triangulation of the input. Duplicates are detected. \code{move(v, p)} moves a vertex (in place if its star stays valid
and locally Delaunay, otherwise by removal, filling the hole with the Delaunay triangulation of its link, and reinsertion).

\paragraph*{Periodic Voronoi tessellation with MPI.} \cls{PeriodicVoronoi}$\langle D\rangle$ is collective: each task gives its
points (any distribution: slabs, pencils, Peano--Hilbert ranges, random, empty tasks) and gets the Voronoi cells of its points
(volume, centroid) and every face touching them (area, centroid, normal, the neighbour's global id, task and index, the periodic
shift). Each task triangulates its points with \emph{ghosts}, copies of other tasks' points and periodic images: an initial
buffer of two mean spacings around the bounding box of its points, then rounds of \emph{certification}, in which the
circumsphere of every simplex with a local vertex is searched on every task (and image) whose bounding box it meets; the points
found are inserted, and a simplex that was searched and survives is final. This needs no assumption on how the points are
distributed and is exact. Coordinates are quantized to $2^{-40}$ of the box so that images are exact translates, perturbation keys
are (global id, image shift), and circumcentres, face rings and sums are computed in an order fixed by the keys, so the cells and
faces are \emph{bitwise identical for any number of tasks and any distribution}, and a face shared by two tasks is identical on
both (a flux can be computed once). The geometry is computed with OpenMP.

\paragraph*{Moving points.} \code{update(new_points)} gives exactly (bitwise) what a new tessellation of the new positions gives,
reusing the ghosts and the triangulation: the ghosts follow their points, points may cross the box and other tasks' bounding boxes
(positions carry a wrap count), and the simplices are certified again. \code{update_method = Automatic} chooses from sampled costs
between moving the vertices and retriangulating; it pays off when few points move (with 1\% of $2\times10^5$ points moving the
triangulation takes 0.39\,s instead of 1.71\,s); when all points move by more than $\sim1\%$ of the spacing an update costs about
as much as a build. The geometry is still recomputed in full.

\begin{cppcode}[Periodic Voronoi tessellation with MPI]
#include <FML/Triangulation/PeriodicVoronoi.h>
using namespace FML::DELAUNAY;
std::vector<std::array<double, 3>> points = ...;  // this task's points in [0, L)^3
PeriodicVoronoi<3> voronoi(points, L);            // collective
for (size_t i = 0; i < voronoi.ncells(); i++) {
    const auto & cell = voronoi.cells()[i];       // cell.id, cell.volume, cell.centroid
    auto [begin, end] = voronoi.cell_faces(i);    // codes of the faces of cell i
    // faces()[code >= 0 ? code : -code - 1]: area, centroid, normal, neighbour, shift, tasks
}
voronoi.update(new_points);                       // = a new build, reusing ghosts and triangulation
\end{cppcode}

\paragraph*{In the library.} \cls{PeriodicVoronoiTessellation}$\langle D\rangle$ (\file{Triangulation/VoronoiTessellation.h})
gives the tessellation in the form the library's users need: the Voronoi volume and the neighbours (Delaunay edges) of every local
particle, with neighbours on other tasks as boundary particles whose values are fetched from their owners. Unlike
\cls{MPIPeriodicDelaunay} there is no buffer or guard points to choose, the particles can be on any domain, and every local cell
is exact. \code{WatershedOnGraph} (\file{Triangulation/Watershed.h}) is the watershed on the neighbour graph of either
tessellation (steepest descent to a local minimum, ties broken by position, links between basins at the lowest connecting
value), \code{WatershedDensityNative} the density watershed on the native one. ZOBOV (\code{ZobovOptions::tessellation =
Native}) and the ORIGAMI halos (\code{OrigamiHaloOptions::tessellation}, \cref{sec:web-origami}) use it by default in builds
without CGAL, as do the Analysis library (\code{zobov_tessellation}) and SnapshotAnalyzer.

\begin{validation}[Validation: native triangulation]
\file{test_delaunay} (69 checks): the predicates equal exact 128-bit integer determinants for $4\times10^5$ small-integer cases
(thousands exactly degenerate) and the filter agrees with exact evaluation on $10^5$ nearly degenerate configurations;
triangulations of random, clustered, lattice and duplicated points are valid and Delaunay (brute force: no vertex inside any
circumsphere); periodic volumes sum to the box to $10^{-15}$, cells close to $10^{-13}$, divergence-theorem volumes agree to
$10^{-12}$, lattices give cubes, results are translation and thread invariant; \code{move}, \code{retriangulate} and
\code{update} equal fresh triangulations. \file{test_delaunay_mpi} (69 checks on 1--4 tasks): slab, pencil, Peano--Hilbert, random
and last-task distributions give cells and faces bitwise identical to all points on task 0, also after updates with points
crossing task and box boundaries. \file{compare_cgal}: the Delaunay edges (with periodic offsets) are identical to CGAL~5.6's
periodic triangulations for $10^5$ and $10^6$ random and clustered points in 2D and 3D, and the triangulation is 1.5--4 times
faster (3D, $10^6$ points: 7.3\,s against 18.0\,s on one core). \file{tests/test_watershed_native} and
\file{tests/test_origami_halos}: native Voronoi volumes, basins, links, ZOBOV catalogues and ORIGAMI halos are independent of how
the particles are spread over the tasks and identical to CGAL's in 2D and 3D. \file{compare_library}: volumes, neighbour graph
and DTFE vertex densities against \cls{MPIPeriodicDelaunay}. \file{tests/test_dtfe_native}: native DTFE grids independent of the particle distribution (slabs, round robin, pencils) and equal
to CGAL's to $5\times10^{-12}$. On a $128^3$ COLA snapshot (200\,Mpc$/h$, $z=0$) the native tessellation gives the same ZOBOV voids
(8129), ORIGAMI halos (5148, every particle in the same halo) and DTFE grids (to $10^{-10}$) as CGAL on 1, 2 and 4 tasks, 1.2--3 times
faster, with a peak memory of 0.88 against 1.10\,GB per task on 4 tasks (2.66 against 2.06\,GB on one task).
\end{validation}

\begin{planned}[Planned]
Next for the moving mesh: recomputing only the geometry around moved points, and load balancing.
\end{planned}

\section{Periodic Delaunay tessellation with CGAL and the DTFE}\label{sec:voids-dtfe}

The DTFE (\cls{DTFE}, \code{DTFETessellation}) uses \fml{}'s own tessellation by default: the particles can be on any task, and the
triangulation is made exact over the grid cells each task fills; \code{DTFETessellation::CGAL} uses \cls{MPIPeriodicDelaunay} instead.
\cls{MPIPeriodicDelaunay} builds the 3D Delaunay tessellation of particles spread over the tasks with CGAL: each task triangulates its
particles plus a buffer from the neighbours (on the domain of the particles, slabs or pencils: the one given to the constructor or the
default particle domain; \cref{sec:par-communication}), and cells whose tetrahedra might be affected by the
buffer edge are flagged as unreliable. Voronoi volumes are the volumes of the contiguous Voronoi cells.

The Delaunay tessellation field estimator \cite{Schaap2000,Cautun2011} assigns each particle the density
\begin{equation}
\rho_i=\frac{(D+1)\,m_i}{W_i},
\end{equation}
with $W_i$ the volume of all tetrahedra sharing the particle, and interpolates any quantity known at the particles linearly inside each
tetrahedron, so gradients are constant per tetrahedron and the density integrates exactly to the total mass. Grid values are the field at the grid
point or the average over the cell ($n_{\rm sub}^3$ samples, volume- or mass-weighted). For velocities the DTFE gives the velocity, divergence,
vorticity and shear grids, the input of the V-web (\cref{sec:web-tweb}). Densities are in units of the mean; positions in $[0,1)$.

\begin{cppcode}[DTFE density and velocity fields]
#include <FML/Triangulation/DTFE.h>
FML::TRIANGULATION::DTFE<Particle> dtfe;
dtfe.create(part.get_particles_ptr(), part.get_npart());
FML::GRID::FFTWGrid<3> density;
dtfe.density_to_grid(Nmesh, density);                    // 1 + delta
FML::TRIANGULATION::DTFEVelocityGrids vel;
dtfe.velocity_to_grid(part.get_particles_ptr(), Nmesh, vel);   // divergence, vorticity, shear
\end{cppcode}

\section{Void finders}\label{sec:voids-finders}

All finders take particles with wrapped positions owned by their slab and return a catalogue of \cls{Void}s that is the same on every task: centre
(circumcentre of the minimum and its three lowest-density mutually adjacent neighbours for ZOBOV, the minimum for voxels), barycentre, volume,
effective radius $R=(3V/4\pi)^{1/3}$, minimum and mean density, the significance (the density at which the void links to a neighbour divided by its
minimum, ZOBOV's density ratio), membership counts, the merging hierarchy, and $\lambda=\delta_{\rm avg}R^{1.2}$ \cite{Nadathur2019}. Each can also
find overdense ``clusters'' by running on $1/\rho$.
\begin{itemize}
\item \emph{Spherical} \cite{Padilla2005}: underdense candidates on a density grid are grown to the largest sphere with mean density below
      $(1+\Delta)\bar\rho$, optionally recentred, and overlapping spheres removed in descending radius (\code{max_overlap}). Each task gets the
      particles within the search radius from every slab that needs them.
\item \emph{Voxel} (Revolver \cite{Nadathur2019}): a watershed of a smoothed CIC density grid; \code{VoxelWatershed} works on any density grid, e.g.\
      the DTFE's.
\item \emph{ZOBOV} \cite{Neyrinck2008}: the watershed of the Voronoi densities of the particles; zones are the basins of local minima.
      \code{ZobovOptions::tessellation} selects the native tessellation (\code{Native}, the default) or CGAL's
      (\code{CGAL}); both give the same catalogue.
\item \emph{Zone merging}: \code{None} (every zone is a void, as Revolver), \code{Threshold} (merge through links below $0.2\bar\rho$, as VIDE
      \cite{Sutter2015}) or \code{Full} (the ZOBOV hierarchy with parents).
\end{itemize}
The default grid (\code{Nmesh = 0}) is Revolver's, $\lfloor2(4\pi N/3)^{1/3}\rfloor$ rounded to a multiple of 8 (32 above 256), independent of the
number of tasks; for a given grid the spherical and voxel catalogues are identical for any number of tasks, and ZOBOV's does not depend on a grid.

\begin{validation}[Validation: tessellations and voids]
\file{tests/test_delaunay_domain.cpp}: on slabs and every pencil layout the same Voronoi volumes, no unreliable cells, and the same watershed basins;
\file{tests/test_dtfe_pencil.cpp}: DTFE densities, gradients and velocity fields on pencils equal those on slabs. \file{test_spherical.cpp}: exact ghost
IDs with tiny messages, periodic and slab boundaries, buffers wider than a slab, empty and non-uniform slabs, and a brute-force catalogue reference (radii,
centres, masses, members, overlap removal); \file{tests/test_voids_pencil.cpp}: voxel and spherical catalogues on pencils equal those on slabs;
\file{tests/test_watershed_native.cpp}: ZOBOV on the native tessellation independent of the particle distribution and identical to CGAL's. The
voxel and merging code has no unit test of its own; it was checked against an independent watershed implementation.
\end{validation}
